\section{Related Patterns}
TODO:

Maybe these are candidates for new patterns

\begin{itemize}
    \item \textbf{Accessible by Design:} Designing teaching materials and tools with accessibility as a first-class concern rather than as an afterthought is a core principle of inclusive education and universal design \cite{w3c2023,burgstahler2015}. The AI tutor pattern complements this approach by enabling accessibility even when legacy materials are not fully accessible.
    \item \textbf{Human-in-the-Loop Accessibility:} Instructors remain responsible for structuring and contextualizing materials, while the AI performs repetitive transformation and explanation tasks. This aligns with human-centered and human-in-the-loop AI concepts, where automation supports but does not replace human expertise \cite{shneiderman2020,ieee2019}.
    \item \textbf{Personal Learning Tutor:} The AI tutor acts as an individualized learning assistant that adapts explanations to the learner’s pace and prior knowledge. Such adaptive tutoring systems are a well-established application area of AI in education \cite{holmes2019}.
    \item \textbf{Progressive Disclosure of Complexity:} Students can request explanations at different levels of detail, supporting incremental learning and multiple representations of knowledge, which are central principles of Universal Design for Learning \cite{burgstahler2015,al-azawei2016}.
    \item \textbf{Single Source of Truth:} Teaching materials remain unchanged for the general audience while accessibility is generated dynamically through the AI. This reduces parallel document maintenance and supports sustainable accessibility practices \cite{lazar2015}.
\end{itemize}
